\chapter[Many-sorted presentation]{Many-sorted presentation of higher-order categories}
\label{ch:many_sorted}

In the many-sorted presentation of higher-order categories, morphisms at each dimension live in
separate sets, related by source, target, and identity maps across different dimensions. Composition
at each pair of indices is defined as a partial operation on morphisms of the higher dimension, only
when appropriate source-target matching conditions are satisfied. In this chapter, we present the
basic definitions and axioms of many-sorted categories and many-sorted functors between them,
concluding with the definition of the category of many-sorted categories itself.

\section{Many-sorted categories}

We start by defining the necessary data for defining a many-sorted category.

\begin{definition}[Many-sorted category structure]
  \label{def:ManySortedCategoryStruct}
  \lean{HigherCategoryTheory.ManySorted.CategoryStruct}
  \leanok
  Let $(C_k)_{k \in \Index}$ be a family of sets whose elements will be called \emph{$k$-morphisms}.
  The required data for defining a many-sorted $\Index$-category structure on $(C_k)_{k \in \Index}$
  consists of, for each pair of indices $k, j \in \Index$ with $j < k$:

  \begin{enumerate}
    \item \label{ax:ms_sc} A map $\Sc^{(k,j)} \colon C_k \to C_j$, called the \emph{$(k,j)$-source}.
    \item \label{ax:ms_tg} A map $\Tg^{(k,j)} \colon C_k \to C_j$, called the \emph{$(k,j)$-target}.
    \item \label{ax:ms_idm} A map $\Idm^{(k,j)} \colon C_j \to C_k$, called the
      \emph{$(k,j)$-identity}.
    \item \label{ax:ms_comp} A partial binary operation $\circ^{(k,j)} \colon C_k \times C_k
      \rightharpoonup C_k$, called the \emph{$(k,j)$-composition}.
  \end{enumerate}

  The composition $g \circ^{(k,j)} f$ is defined if and only if $\Sc^{(k,j)} (g) = \Tg^{(k,j)} (f)$,
  in which case we say that $g$ and $f$ are \emph{$(k,j)$-composable}.

  We will represent the above data as a tuple $\cat{C} = ((C_k)_{k \in \Index}, (\Sc^{(k,j)},
  \Tg^{(k,j)}, \Idm^{(k,j)}, \circ^{(k,j)})_{k, j \in \Index, \, j < k})$.
\end{definition}

We split the axioms of many-sorted categories into two parts: first, those concerning only one pair
of indices at a time, and then those concerning the interaction between different pairs. This
division is made by first defining the notion of pre-many-sorted category:

\begin{definition}[Pre-many-sorted category]
  \label{def:PreManySortedCategory}
  \uses{def:ManySortedCategoryStruct}
  \lean{HigherCategoryTheory.ManySorted.PreCategory}
  \leanok
  A \emph{pre-many-sorted $\Index$-category} is a tuple
  \[
    \cat{C} = ((C_k)_{k \in \Index}, (\Sc^{(k,j)}, \Tg^{(k,j)}, \Idm^{(k,j)},
    \circ^{(k,j)})_{k, j \in \Index, \, j < k}),
  \]
  satisfying, for each pair of indices $k, j \in \Index$ with $j < k$, the following axioms:

  \begin{enumerate}
    \item \label{ax:ms_sc_tg_comp} For all $f, g \in C_k$, if $g$ and $f$ are $(k,j)$-composable,
      then
      \begin{align*}
        \Sc^{(k,j)} (g \circ^{(k,j)} f) = \Sc^{(k,j)} (f), \qquad &
        \Tg^{(k,j)} (g \circ^{(k,j)} f) = \Tg^{(k,j)} (g).
      \end{align*}
    \item \label{ax:ms_sc_tg_idm} For all $f \in C_j$,
      \[
        \Sc^{(k,j)} (\Idm^{(k,j)} (f)) = f, \qquad \Tg^{(k,j)} (\Idm^{(k,j)} (f)) = f.
      \]
    \item \label{ax:ms_id_laws} For all $f \in C_k$, the morphisms $f$ and $\Idm^{(k,j)}
      (\Sc^{(k,j)} (f))$ are $(k,j)$-composable and
      \[
        f \circ^{(k,j)} \Idm^{(k,j)} (\Sc^{(k,j)} (f)) = f.
      \]
      Similarly, $\Idm^{(k,j)} (\Tg^{(k,j)} (f))$ and $f$ are $(k,j)$-composable and
      \[
        \Idm^{(k,j)} (\Tg^{(k,j)} (f)) \circ^{(k,j)} f = f.
      \]
    \item \label{ax:ms_assoc} For all $f, g, h \in C_k$, if $h$ and $g$ are $(k,j)$-composable, and
      $g$ and $f$ are $(k,j)$-composable, then $h \circ^{(k,j)} g$ and $f$ and $h$ and $g
      \circ^{(k,j)} f$ are $(k,j)$-composable, and
      \[
        (h \circ^{(k,j)} g) \circ^{(k,j)} f = h \circ^{(k,j)} (g \circ^{(k,j)} f).
      \]
  \end{enumerate}

\end{definition}

An immediate consequence of the source/target axioms for identities is that the identity map at each
pair of dimensions is injective:

\begin{proposition}
  \label{prop:PreManySortedCategory.injetive_idm}
  \uses{def:PreManySortedCategory}
  \lean{HigherCategoryTheory.ManySorted.PreCategory.injetive_idm}
  \leanok
  Let $\cat{C}$ be a pre-many-sorted $\Index$-category and let $j < k$ in $\Index$. Then the
  identity map $\Idm^{(k,j)} \colon C_j \to C_k$ is injective.
\end{proposition}

\begin{proof}
  \leanok
  Let $f, g \in C_j$ with $\Idm^{(k,j)}(f) = \Idm^{(k,j)}(g)$. Then
  \[
    f = \Sc^{(k,j)}(\Idm^{(k,j)}(f)) = \Sc^{(k,j)}(\Idm^{(k,j)}(g)) = g,
  \]
  where both equalities use Axiom~\refaxiom{def:PreManySortedCategory}{ax:ms_sc_tg_idm}.
\end{proof}

We now define a many-sorted category structure by extending the above definition with the axioms
concerning the interaction between different pairs of indices:

\begin{definition}[Many-sorted category]
  \label{def:ManySortedCategory}
  \uses{def:PreManySortedCategory}
  \lean{HigherCategoryTheory.ManySorted.Category}
  \leanok
  A \emph{many-sorted $\Index$-category} is a pre-many-sorted $\Index$-category $\cat{C} = ((C_k)_{k
  \in \Index}, (\Sc^{(k,j)}, \Tg^{(k,j)}, \Idm^{(k,j)}, \circ^{(k,j)})_{k, j \in \Index, \, j < k})$
  satisfying, in addition to the axioms of pre-many-sorted categories, the following
  cross-dimensional axioms, for all indices $k, j, i \in \Index$ with $i < j < k$:

  \begin{enumerate}
    \item \label{ax:ms_cross_glob} For all $f \in C_k$,
      \begin{align*}
        \Sc^{(j,i)} (\Sc^{(k,j)} (f)) = \Sc^{(k,i)} (f), \qquad &
        \Sc^{(j,i)} (\Tg^{(k,j)} (f)) = \Sc^{(k,i)} (f), \\
        \Tg^{(j,i)} (\Tg^{(k,j)} (f)) = \Tg^{(k,i)} (f), \qquad &
        \Tg^{(j,i)} (\Sc^{(k,j)} (f)) = \Tg^{(k,i)} (f).
      \end{align*}
    \item \label{ax:ms_cross_sc_tg_comp} For all $f, g \in C_k$, if $g$ and $f$ are
      $(k,i)$-composable, then $\Sc^{(k,j)} (g)$ and $\Sc^{(k,j)} (f)$ and $\Tg^{(k,j)} (g)$ and
      $\Tg^{(k,j)} (f)$ are also $(j,i)$-composable, and
      \begin{align*}
        \Sc^{(k,j)} (g \circ^{(k,i)} f) = \Sc^{(k,j)} (g) \circ^{(j,i)} \Sc^{(k,j)} (f), \\
        \Tg^{(k,j)} (g \circ^{(k,i)} f) = \Tg^{(k,j)} (g) \circ^{(j,i)} \Tg^{(k,j)} (f).
      \end{align*}
    \item \label{ax:ms_idm_idm} For all $f \in C_i$,
      \[
        \Idm^{(k,j)} (\Idm^{(j,i)} (f)) = \Idm^{(k,i)} (f).
      \]
    \item \label{ax:ms_idm_comp} For all $f, g \in C_j$, if $g$ and $f$ are $(j,i)$-composable, then
      $\Idm^{(k,j)} (g)$ and $\Idm^{(k,j)} (f)$ are $(k,i)$-composable, and
      \[
        \Idm^{(k,j)} (g \circ^{(j,i)} f) = \Idm^{(k,j)} (g) \circ^{(k,i)} \Idm^{(k,j)} (f).
      \]
    \item \label{ax:ms_interchange} For all $f_1, f_2, g_1, g_2 \in C_k$, suppose that:
      \begin{itemize}
        \item $g_2$ and $g_1$ are $(k,j)$-composable,
        \item $f_2$ and $f_1$ are $(k,j)$-composable,
        \item $g_2$ and $f_2$ are $(k,i)$-composable, and
        \item $g_1$ and $f_1$ are $(k,i)$-composable.
      \end{itemize}
      Then, the morphisms $(g_2 \circ^{(k,i)} f_2)$ and $(g_1 \circ^{(k,i)} f_1)$ are
      $(k,j)$-composable, and the morphisms $(g_2 \circ^{(k,j)} g_1)$ and $(f_2 \circ^{(k,j)} f_1)$
      are $(k,i)$-composable, and
      \[
        (g_2 \circ^{(k,i)} f_2) \circ^{(k,j)} (g_1 \circ^{(k,i)} f_1) =
        (g_2 \circ^{(k,j)} g_1) \circ^{(k,i)} (f_2 \circ^{(k,j)} f_1).
      \]
      That is, composing first at dimension $j$ and then at dimension $i$ yields the same result as
      composing first at dimension $i$ and then at dimension $j$.
  \end{enumerate}

  \begin{definition}[Many-sorted finite category]
    \label{def:ManySortedNCategory}
    \uses{def:ManySortedCategory}
    \lean{HigherCategoryTheory.ManySorted.NCategory}
    \leanok
    Given a natural number $n \in \NN$, a \emph{many-sorted $n$-category} is a many-sorted category
    indexed by the finite set $n + 1 = \set{0, 1, \ldots, n}$.
  \end{definition}

  \begin{remark}
    \label{rem:PreManySortedCategoryLift}
    \uses{def:ManySortedNCategory}
    \lean{HigherCategoryTheory.ManySorted.PreCategory.lift}
    Any pre-many-sorted $1$-category is, in fact, a many-sorted $1$-category. Since the index
    set $1 + 1 = \set{0, 1}$ has exactly two elements, there are no triples of distinct indices $i <
    j < k$, so all cross-dimensional axioms are vacuously satisfied.
  \end{remark}

  \begin{definition}[Many-sorted omega category]
    \label{def:ManySortedOmegaCategory}
    \uses{def:ManySortedCategory}
    \lean{HigherCategoryTheory.ManySorted.OmegaCategory}
    \leanok
    A \emph{many-sorted $\omega$-category} is a many-sorted category indexed by the set of natural
    numbers $\NN$.
  \end{definition}
\end{definition}

\section{Many-sorted functors}

In the many-sorted setting, a functor is a family of maps, one for each dimension, that collectively
preserve sources, targets, identities, and composition at each pair of dimensions.

\begin{definition}[Many-sorted functor]
  \label{def:ManySortedFunctor}
  \uses{def:ManySortedCategory}
  \lean{HigherCategoryTheory.ManySorted.Functor}
  \leanok
  Let $\cat{C} = ((C_k)_{k \in \Index}, (\Sc^{(k,j)}, \Tg^{(k,j)}, \Idm^{(k,j)}, \circ^{(k,j)})_{k,
  j \in \Index, \, j < k})$ and $\cat{D} = ((D_k)_{k \in \Index}, (\Sc^{(k,j)}, \Tg^{(k,j)},
  \Idm^{(k,j)}, \circ^{(k,j)})_{k, j \in \Index, \, j < k})$ be many-sorted $\Index$-categories. A
  \emph{many-sorted $\Index$-functor} from $\cat{C}$ to $\cat{D}$ is a family of maps $F = (F_k
  \colon C_k \to D_k)_{k \in \Index}$, written $F \colon \cat{C} \to \cat{D}$, satisfying, for each
  pair of indices $k, j \in \Index$ with $j < k$, the following preservation properties:

  \begin{enumerate}
    \item \label{ax:ms_fun_sc_tg} For all $f \in C_k$,
      \[
        F_j (\Sc^{(k,j)} (f)) = \Sc^{(k,j)} (F_k (f)), \qquad
        F_j (\Tg^{(k,j)} (f)) = \Tg^{(k,j)} (F_k (f)).
      \]
    \item \label{ax:ms_fun_idm} For all $f \in C_j$,
      \[
        F_k (\Idm^{(k,j)} (f)) = \Idm^{(k,j)} (F_j (f)).
      \]
    \item \label{ax:ms_fun_comp} For all $f, g \in C_k$, if $g$ and $f$ are
      $(k,j)$-composable, then $F_k (g)$ and $F_k (f)$ are also $(k,j)$-composable, and
      \[
        F_k (g \circ^{(k,j)} f) = F_k (g) \circ^{(k,j)} F_k (f).
      \]
  \end{enumerate}
\end{definition}

We can define composition of many-sorted functors componentwise, and the identity many-sorted
functor as the family of identity functions.

\begin{definition}[Composition of many-sorted functors]
  \label{def:ManySortedFunctor_comp}
  \uses{def:ManySortedFunctor}
  \lean{HigherCategoryTheory.ManySorted.Functor.comp}
  \leanok
  Let $\cat{C}, \cat{D}, \cat{E}$ be many-sorted $\Index$-categories, and let $F = (F_k)_{k \in
  \Index} \colon \cat{C} \to \cat{D}$ and $G = (G_k)_{k \in \Index} \colon \cat{D} \to \cat{E}$ be
  many-sorted $\Index$-functors. Then, the family $G \circ F = (G_k \circ F_k)_{k \in \Index} \colon
  \cat{C} \to \cat{E}$ is also a many-sorted $\Index$-functor.
\end{definition}

\begin{proof}
  \leanok
  Let $k, j \in \Index$ with $j < k$. For all $f \in C_k$, since $F$ and $G$ are many-sorted
  functors, we have that
  \[
    \begin{split}
      (G_j \circ F_j) (\Sc^{(k,j)} (f))
      & = G_j (F_j (\Sc^{(k,j)} (f))) \\
      & = G_j (\Sc^{(k,j)} (F_k (f))) \\
      & = \Sc^{(k,j)} (G_k (F_k (f))) \\
      & = \Sc^{(k,j)} ((G_k \circ F_k) (f)),
    \end{split}
  \]
  and similarly,
  \[
    \begin{split}
      (G_j \circ F_j) (\Tg^{(k,j)} (f))
      & = G_j (F_j (\Tg^{(k,j)} (f))) \\
      & = G_j (\Tg^{(k,j)} (F_k (f))) \\
      & = \Tg^{(k,j)} (G_k (F_k (f))) \\
      & = \Tg^{(k,j)} ((G_k \circ F_k) (f)),
    \end{split}
  \]
  that is, $G \circ F$ preserves sources and targets at pair $(k,j)$.

  For all $f \in C_j$, we have that
  \[
    \begin{split}
      (G_k \circ F_k) (\Idm^{(k,j)} (f))
      & = G_k (F_k (\Idm^{(k,j)} (f))) \\
      & = G_k (\Idm^{(k,j)} (F_j (f))) \\
      & = \Idm^{(k,j)} (G_j (F_j (f))) \\
      & = \Idm^{(k,j)} ((G_j \circ F_j) (f)),
    \end{split}
  \]
  so $G \circ F$ preserves identities at pair $(k,j)$.

  Now, let $f, g \in C_k$ be $(k,j)$-composable. We have that
  \[
    \begin{split}
      (G_k \circ F_k) (g \circ^{(k,j)} f)
      & = G_k (F_k (g \circ^{(k,j)} f)) \\
      & = G_k (F_k (g) \circ^{(k,j)} F_k (f)) \\
      & = G_k (F_k (g)) \circ^{(k,j)} G_k (F_k (f)) \\
      & = (G_k \circ F_k) (g) \circ^{(k,j)} (G_k \circ F_k) (f),
    \end{split}
  \]
  so $G \circ F$ preserves composition at pair $(k,j)$.
\end{proof}

\begin{definition}[Identity many-sorted functor]
  \label{def:ManySortedFunctor_id}
  \uses{def:ManySortedFunctor}
  \lean{HigherCategoryTheory.ManySorted.Functor.id}
  \leanok
  Let $\cat{C}$ be a many-sorted $\Index$-category. Then, the family of identity functions
  $\id_{\cat{C}} = (\id_{C_k})_{k \in \Index}$, written $\id_{\cat{C}} \colon \cat{C} \to \cat{C}$,
  is a many-sorted $\Index$-functor from $\cat{C}$ to itself.
\end{definition}

\begin{proof}
  \leanok
  It trivially follows from reflexivity.
\end{proof}

\begin{definition}[Many-sorted finite functor]
  \label{def:ManySortedNFunctor}
  \uses{def:ManySortedNCategory, def:ManySortedFunctor}
  \lean{HigherCategoryTheory.ManySorted.NFunctor}
  \leanok
  Given a natural number $n \in \NN$, a \emph{many-sorted $n$-functor} is a many-sorted functor
  between many-sorted $n$-categories.
\end{definition}

\begin{definition}[Many-sorted omega functor]
  \label{def:ManySortedOmegaFunctor}
  \uses{def:ManySortedOmegaCategory, def:ManySortedFunctor}
  \lean{HigherCategoryTheory.ManySorted.OmegaFunctor}
  \leanok
  A \emph{many-sorted $\omega$-functor} is a many-sorted functor between many-sorted
  $\omega$-categories.
\end{definition}

\section{The category of many-sorted categories}

Lastly, we define the category whose objects are many-sorted categories and whose morphisms are
many-sorted functors.

\begin{definition}[Category of many-sorted categories]
  \label{def:ManySortedCategory_category}
  \uses{
    def:ManySortedCategory,
    def:ManySortedFunctor,
    def:ManySortedFunctor_comp,
    def:ManySortedFunctor_id
  }
  \lean{HigherCategoryTheory.ManySorted.Cat.category}
  \leanok
  The \emph{category of many-sorted $\Index$-categories}, denoted by $\Index\CatMS$, is defined as
  follows:

  \begin{enumerate}
    \item \label{ax:ms_cat_obj} Its objects are many-sorted $\Index$-categories.
    \item \label{ax:ms_cat_mor} Its morphisms are many-sorted $\Index$-functors between many-sorted
      $\Index$-categories.
    \item \label{ax:ms_cat_comp} Composition of morphisms is given by composition of many-sorted
      functors.
    \item \label{ax:ms_cat_id} The identity morphism for each object is given by the identity
      many-sorted functor.
  \end{enumerate}
\end{definition}

\begin{proof}
  \leanok
  Composition and the identity functor are well defined by
  Definitions~\ref{def:ManySortedFunctor_comp} and~\ref{def:ManySortedFunctor_id}, respectively. The
  required category laws trivially follow from the corresponding properties of function composition.
\end{proof}

\begin{definition}[Category of many-sorted finite categories]
  \label{def:ManySortedNCategory_category}
  \uses{
    def:ManySortedNCategory,
    def:ManySortedNFunctor,
    def:ManySortedFunctor_comp,
    def:ManySortedFunctor_id
  }
  \lean{HigherCategoryTheory.ManySorted.NCat.category}
  \leanok
  Given a natural number $n \in \NN$, the \emph{category of many-sorted $n$-categories}, denoted by
  $n\CatMS$, is the particular case of Definition~\ref{def:ManySortedCategory_category} where
  $\Index$ is the finite set $n + 1 = \set{0, 1, \ldots, n}$. Its objects are many-sorted
  $n$-categories and its morphisms are many-sorted $n$-functors.
\end{definition}

\begin{definition}[Category of many-sorted omega categories]
  \label{def:ManySortedOmegaCategory_category}
  \uses{
    def:ManySortedOmegaCategory,
    def:ManySortedOmegaFunctor,
    def:ManySortedFunctor_comp,
    def:ManySortedFunctor_id
  }
  \lean{HigherCategoryTheory.ManySorted.OmegaCat.category}
  \leanok
  The \emph{category of many-sorted $\omega$-categories}, denoted by $\omega\CatMS$, is the
  particular case of Definition~\ref{def:ManySortedCategory_category} where $\Index$ is the set of
  natural numbers $\NN$. Its objects are many-sorted $\omega$-categories and its morphisms are
  many-sorted $\omega$-functors.
\end{definition}
